\section{Probleme und Herausforderungen}
In diesem Abschnitt beschreibe ich die wichtigsten technischen Schwierigkeiten, die im Verlauf der Implementierung aufgetreten sind.
Die Darstellung ist bewusst problemorientiert, da die Architektur und die Algorithmik bereits in Kapitel~4 bzw.\ Kapitel~2 erläutert wurden.

\subsection{Fehlende explizite Chance Nodes für öffentliche Karten}
Ein zentrales Problem entstand dadurch, dass öffentliche Karten in der Game Environment automatisch am Ende einer Setzrunde ausgeteilt werden.
Dadurch existieren im Spielbaum keine expliziten public chance nodes, die den Kartendeal als eigene Zustandsübergänge repräsentieren.
Für viele Komponenten ist diese Unterscheidung jedoch wichtig, weil Chance-Ereignisse sonst nicht sauber modelliert und wiederverwendet werden können.

\subsection{Konsequenzen für Public State Tree und Best Response}
Beim Aufbau des Public State Trees zeigte sich, dass der Best-Response-Ansatz explizite Chance-Knoten erwartet, um Zustände korrekt zu gruppieren und Reach-Probabilities konsistent zu berechnen.
Um das in der bestehenden Game Environment dennoch abbilden zu können, war zusätzlicher Code notwendig, der den öffentlichen Verlauf rekonstruiert und die öffentlichen Karten an geeigneten Stellen „nachträglich“ behandelt.
Diese Lösung machte die Evaluation funktionsfähig, erhöhte aber die Komplexität und erschwerte spätere strukturelle Änderungen.

\subsection{Sampling-Solver: Ziel nicht erreicht}
Zusätzlich habe ich Sampling-basierte Solver implementiert, um die in der Literatur beschriebenen Skalierungsvorteile praktisch zu nutzen.
In meinem Implementierungsstand habe ich diese Solver jedoch nicht zuverlässig zum gewünschten Konvergenzverhalten gebracht.
Rückblickend ist es plausibel, dass dies prinzipiell lösbar ist, jedoch eine wesentlich sorgfältigere Behandlung der Chance-Ereignisse und des Deal-Samplings erfordert hätte.
Da diese Arbeiten den zeitlichen Rahmen gesprengt hätten, wurden die Sampling-Solver im weiteren Projektverlauf nicht mehr als Kernvergleichsbaseline betrachtet.

\subsection{Skalierungsgrenzen am Beispiel großer Hold'em-Varianten}
Mit steigender Spielkomplexität verschieben sich die praktischen Engpässe von „funktioniert grundsätzlich“ hin zu „passt noch in Speicher und Laufzeitbudget“.
Für Heads-Up Limit Hold'em war in meinem Projektstand insbesondere der Zustandsraum so groß, dass Tree-basierte Ansätze in der Praxis nicht mehr einsetzbar waren.
Für solche Spiele fallen damit alle Solver-Varianten weg, die einen vollständigen Game Tree vorab materialisieren müssen.
Gleichzeitig bedeutet das nicht, dass Training in größeren Spielen grundsätzlich unmöglich ist, sondern dass die aktuelle Implementierung ohne weitere Maßnahmen sehr langsam wird.
Als grobe Größenordnung habe ich für Rhode Island Hold'em geschätzt, dass eine einzelne Iteration in meinem aktuellen Implementierungsstand etwa 15 Stunden dauern würde.
Um große Varianten wie Limit Hold'em praktisch bearbeiten zu können, wären daher mehrere weitere Schritte notwendig.
Erstens wäre eine stärkere Optimierung der Implementation erforderlich, da in Python ein erheblicher Overhead entsteht und sich durch effizientere Datenstrukturen oder native Komponenten vermutlich viel Laufzeit einsparen lässt.
Zweitens ist Abstraktion notwendig, um den Zustandsraum so zu reduzieren, dass Training und Evaluation in realistischer Zeit möglich werden.
Drittens wäre Parallelisierung sinnvoll, um unabhängige Teilprobleme (z.B.\ Deals oder Teilbäume) auf mehrere Kerne zu verteilen.
Viertens sind zusätzliche Optimierungsverfahren wie Pruning oder verwandte Beschleunigungen hilfreich, um unnötige Traversierungen zu vermeiden und Rechenzeit gezielt dort zu investieren, wo sie das Ergebnis beeinflusst.

\subsection{Nicht umgesetzte Abstraktion: Hand-strength-squared}
Eine zusätzliche Abstraktionsidee war eine Reduktion über Hand-strength-squared.
Diese Erweiterung wurde im Projektverlauf geplant, konnte jedoch aus Zeitgründen nicht mehr umgesetzt werden.

